\begin{pastexam}[label={px:nchu-113-dsa-17}]{中興資工所甲組 113 學年度 選擇題第 17 題}
（單選；答錯不倒扣）

You are planning a promotional event in a city, where you need to hang banners at various locations. Each
location has a specific time window during which the banner can be hung. Given a list of locations with
their respective time windows represented as pairs (start\_time, end\_time), your task is to maximize the
number of locations where the banners can be hung without overlapping time windows. What is the most
appropriate algorithm for planning the banner placement? (5\%)
\begin{enumerate}[label=(\Alph*)]
\item A greedy algorithm that selects the location with the earliest start time first.
\item A greedy algorithm that selects the location with the shortest available time window first.
\item A greedy algorithm that selects the location with the earliest end time first.
\item A dynamic programming approach that considers all possible combinations of locations.
\end{enumerate}
\end{pastexam}

\begin{pastsoln}
答案是 (C)。

這就是活動選擇問題（資結書第~19.1~節）：選最多個兩兩不重疊的區間。 依結束時間排序，每次選「結束最早、
而且和已選的不重疊」的區間，是最佳解（資結書第~19.1.5~節的交換論證：任何最佳解的第一個區間都可以換成
結束最早的那個，不會變少）。 排序 $O(n \log n)$，掃一遍 $O(n)$。

另外三個：
\begin{itemize}
\item (A) 最早開始：反例 $[0, 10)$、$[1, 2)$、$[3, 4)$。 先選 $[0, 10)$，後兩個都和它重疊，只得 $1$ 個；最佳是
$[1, 2)$、$[3, 4)$ 兩個。
\item (B) 最短區間：反例 $[0, 5)$、$[4, 7)$、$[6, 11)$。 最短的 $[4, 7)$ 和另外兩個都重疊，只得 $1$ 個；最佳是
$[0, 5)$、$[6, 11)$ 兩個。
\item (D) 列舉所有組合是 $2^n$ 種，雖然正確但是指數時間；這題不需要動態規劃（動態規劃用在帶權重的區間排程，
資結書第~18.7.3~節）。
\end{itemize}
\end{pastsoln}
